Choose the repeated color in 5 ways. If its faces are opposite, the four distinct remaining colors occupy a ring, modulo rotations and reflections of that ring, giving $\dfrac{4!}{8}=3$ arrangements. If its faces are adjacent, fix the shared pair: the full cube symmetry group has 48 elements, acts transitively on the 12 adjacent pairs, and therefore has a four-element stabilizer of that pair. Each nonidentity such symmetry moves at least two of the four distinctly colored remaining faces, so the 24 assignments fall into groups of four, giving 6 arrangements. Total $5(3+6)=\boxed{45}$.
\begin{asy}
texpreamble("\usepackage{amsmath,amssymb}");
size(240);
defaultpen(fontsize(9pt));
draw(box((0,0),(2,2)));filldraw((0,2)--(.7,2.7)--(2.7,2.7)--(2,2)--cycle,lightblue);draw((2,0)--(2.7,.7)--(2.7,2.7)--(2,2));label("R",(1.35,2.35));label("R also on bottom",(1,-.6));label("opposite",(1,3.3));
\end{asy}
\begin{asy}
texpreamble("\usepackage{amsmath,amssymb}");
size(240);
defaultpen(fontsize(9pt));
filldraw(box((0,0),(2,2)),lightblue);filldraw((0,2)--(.7,2.7)--(2.7,2.7)--(2,2)--cycle,lightblue);draw((2,0)--(2.7,.7)--(2.7,2.7)--(2,2));label("R",(1,1));label("R",(1.35,2.35));label("adjacent",(1,3.3));
\end{asy}
